\section*{第 \ref{ch:g5:irreversible-changes} 章——不可逆变化}

\begin{solution}{exo:g5:irreversible-changes:1}
熔化巧克力：可以恢复（冷却后会重新凝固）。点燃
一根火柴：无法恢复。溶解食盐：可以恢复（蒸发掉水）。
烤面包：无法恢复。
\end{solution}

\begin{solution}{exo:g5:irreversible-changes:2}
生成新物质的变化。在厨房里：煮鸡蛋、
烤蛋糕（另外：烤面包、切开的苹果变褐）。
\end{solution}

\begin{solution}{exo:g5:irreversible-changes:3}
通过蒸发水分：盐留在碟子里。溶解
是可逆变化。
\end{solution}

\begin{solution}{exo:g5:irreversible-changes:4}
牛奶中出现新的气味，并出现块状固体。
\end{solution}

\begin{solution}{exo:g5:irreversible-changes:5}
左侧的变化：糖溶于水、冰熔化、沙子
与砾石混合。双箭头表示变化可以双向进行，
即可以恢复。
\end{solution}

\begin{solution}{exo:g5:irreversible-changes:6}
油使空气和水无法接触链条的铁，
从而减缓生锈。
\end{solution}

\begin{solution}{exo:g5:irreversible-changes:7}
未经加热，液体中出现气泡。
\end{solution}

\begin{solution}{exo:g5:irreversible-changes:8}
有用的：煮食物、烤面包、燃烧燃气取暖。有害的：
生锈、食物变质、木头腐烂。
\end{solution}

\begin{solution}{exo:g5:irreversible-changes:9}
低温会减缓使牛奶变酸的化学变化，
因此牛奶保存的时间更长。
\end{solution}

\begin{solution}{exo:g5:irreversible-changes:10}
不对。气泡是水蒸气，冷却后会
重新变回水：没有生成新物质，变化可以恢复。
仅有气泡并不能证明出现了新物质。
\end{solution}

\begin{solution}{exo:g5:irreversible-changes:11}
燃烧的蜡烛放出热量和光，并生成新物质：
水蒸气（把冷玻璃杯放在其上方，杯子会起雾）和二氧化碳。
燃烧掉的蜡不能再回来。另一方面，熔化的蜡
冷却后会重新凝固。
\end{solution}
